% Conclusion section for Kinetic AI paper
% All numbers from validated findings F1–F3 (MMD convergence), F6 (auction truthfulness),
% F18 (H1 parity verdict), F20 (scale results), F21 (preference-magnet), F22 (H4 auction),
% F23 (H5 closed-loop), F24 (H6 anytime/certification), F25 (conversion damage).

\section{Conclusion}

This program established five key results and identified sharp open problems where equilibrium-computational framings either succeeded or failed at small-LM scale.

On training dynamics, Magnetic Mirror Descent achieves linear convergence to its regularized fixed point where simultaneous gradient descent-ascent cycles on symmetric zero-sum games ($R^2 = 0.995$ on matching pennies, F1). Periodic reference resets drive Nash convergence to within $10^{-4}$ across all tested games (F3). However, the uniform-anchor fixed point on asymmetric games diverges from the logit-QRE, exposing a limit to the theory when applied outside its proof context (F2). On preference optimization, the magnetic pull to a reference is real but second-order to the preference-optimization gradient across four orders of magnitude of regularization strength ($\tau \in [10^{-3}, 10]$); the pre-registered hypothesis that it would control KL drift below DPO's was not supported (F21). This places magnetic preference anchoring outside its natural home (pretraining stability) and motivates its repositioning in future work.

On model depth, a weight-tied transformer block trained with unrolled anytime supervision on intermediate iterates reaches 99.1\% of a parameter-matched 12-layer explicit baseline's BLiMP score at 121M parameters (mean 0.6637 vs baseline 0.7133; F24, B1 arm). This erases a quality gap that solver-based implicit training had widened from 7\% to 21\% as width scaled, and the improvement was validated across three md5-distinct seeds with tight confidence intervals. The anytime-trained model trains 2.1× faster than implicit-differentiation training and provides a ``think-harder dial'' at inference time: quality degrades gracefully across solver budgets 4, 8, and 12 iterations (F24, B1). The architectural depth-memory advantage is retained: peak activation memory is 23\% lower than the explicit baseline (F20). A separate arm achieves certified fixed-point convergence (rate 1.0 at 4.0 mean iterations) via a solver-aware auxiliary loss, but at a quality cost; the naive combination of anytime supervision and the contraction penalty is refuted by its own pre-registered gate (F24, B2/B4), leaving quality-preserving certification as a precisely formulated open problem.

On decoding, truthful second-price token auctions with model confidence as valuations are empirically exactly truthful across 16,000 observations ($0.0$ regret, 95\% CI $[0.0, 0.0]$, F6). When applied to specialist models trained on disjoint data subsets, the auction achieves 182.8 mixed-domain perplexity, beating the best single specialist (236.8 ppl) by 23\% and uniform logit averaging (207.9 ppl) by 12\% at scoring time (F22). However, the advantage is abolished in closed-loop generation, where context-dependent hand-offs between specialists dominate the per-token selection signal; the evaluation metric itself is shown to be style-dominated, and a style-neutral judge is not available at this scale (F23). The failure was predicted by the adversarial reviewer's insistence on a scoring-time scope, and the follow-up experiment vindicated that distinction; this motivates future work on context-aware bidding and evaluation under style-neutral judges.

The program's negative results carry as much design information as the positive ones. The pre-registered parity gate (≥95\% of explicit baseline under implicit training alone) was formally missed---the best post-LN solver-trained model reached 93.0\% (F18), within noise but below the threshold. This drove the invention of anytime supervision, which then met the quality goal but required training two architecture instances to achieve certification. The closed-loop auction hypothesis (F23) was missed, but its failure exposed a confound in the evaluation (judge bias) that systematic re-analysis then characterized, improving the scientific resolution of the claim.

Future work divides into three directions. First, achieve quality-preserving certification: both halves are now separately solved (parity via anytime, convergence via the auxiliary loss), and their combination is falsifiably broken. Second, convert pretrained models at 1–3B scale using the anytime objective; initial damage curves (F25) show that layers can be selectively tied, and the scale-sensitivity of contraction now drives the next frontier. Third, evaluate auction aggregation under a style-neutral judge and with context-aware bids that key on the input's domain rather than the model's own confidence. These directions are framed as testable experiments with the same pre-registration and audit rigor as the closed empirical program reported here.

